%───────────────────────────────────────────────────────────────────────────────
\chapter{Diagramas}
\label{ch:diagramas}
%───────────────────────────────────────────────────────────────────────────────

Los diagramas en un documento académico no son decoración: son un lenguaje
formal con gramática propia. Un autómata finito es una definición matemática
en forma visual. Un árbol de sintaxis abstracta es la estructura semántica
de un programa. Un diagrama conmutativo es una proposición sobre morfismos.
El template~\cite{tikz2023} trata los diagramas con la misma precisión
que las ecuaciones.

%───────────────────────────────────────────────────────────────────────────────
\section{Filosofía de theming en diagramas}
\label{sec:tikz-filosofia}
%───────────────────────────────────────────────────────────────────────────────

\begin{definicion}{Aliases tikz-* como valores por defecto}
\label{def:tikz-aliases}
El dominio \texttt{tikz-*} de la paleta define valores por defecto para
nodos, aristas y labels genéricos. Son el punto de partida — no el único
color posible. Cuando un nodo tiene significado semántico propio
(estado aceptador, colisión, estado inicial), se sobreescribe con el
alias del dominio correspondiente: \texttt{tikz-accept}, \texttt{fb-error},
\texttt{tikz-initial}.
\end{definicion}

\begin{acadtable}[H]{tab:tikz-aliases}{Aliases semánticos para diagramas TikZ}
  \begin{tabularx}{\textwidth}{l l l X}
    \toprule
    \headerrow
    \thead{Alias} & \thead{Color Ctp} & \thead{Rol} & \thead{Uso} \\
    \midrule
    \texttt{tikz-node-bg}     & Surface0  & Fondo genérico
      & Nodo sin semántica propia \\
    \texttt{tikz-node-border} & Overlay0  & Borde genérico
      & Nodo estándar \\
    \texttt{tikz-edge}        & Overlay1  & Arista
      & Conexión sin dirección semántica \\
    \texttt{tikz-label}       & Subtext0  & Label
      & Anotación secundaria \\
    \texttt{tikz-initial}     & Blue      & Estado inicial
      & Autómatas: punto de entrada \\
    \texttt{tikz-accept}      & Green     & Estado aceptador
      & Autómatas: estado final \\
    \texttt{tikz-reject}      & Red       & Estado rechazador
      & Autómatas / error \\
    \texttt{tikz-highlight}   & Yellow    & Nodo destacado
      & Paso activo en explicaciones \\
    \midrule
    \texttt{fb-error}         & Red       & Error / colisión
      & Condición de error en estructuras \\
    \texttt{eq-known}         & Mauve     & Dato conocido
      & Pesos, etiquetas de aristas con valor fijo \\
    \texttt{env-phys}         & Sapphire  & Transformación
      & Cajas de función / operación \\
    \bottomrule
  \end{tabularx}
\end{acadtable}

La regla es consistente con el resto del template: \textbf{nunca} usar
\texttt{Ctp*} directamente en archivos de contenido. Cuando ningún alias
existente cubre el rol semántico necesario, se añade uno nuevo en
\texttt{themes/catppuccin-palette.tex} bajo el dominio \texttt{tikz-*}.

%───────────────────────────────────────────────────────────────────────────────
\section{Autómatas finitos}
\label{sec:automatas}
%───────────────────────────────────────────────────────────────────────────────

\begin{definicion}{Estilo base para autómatas}
\label{def:automata-style}
El estilo canónico para autómatas en el template usa la librería
\texttt{automata} de TikZ con los aliases semánticos del dominio
\texttt{tikz-*}. El estado inicial usa \texttt{tikz-initial} (Blue),
los estados finales \texttt{tikz-accept} (Green), y las aristas
\texttt{tikz-edge} (Overlay1). Las etiquetas de transición aplican
coloreado semántico de ecuaciones donde corresponda.
\end{definicion}

\begin{lstlisting}[style=latex, frame=none, numbers=none]
\tikzset{
  state/.style = {
    circle,
    draw  = tikz-node-border,
    fill  = background,
    text  = maintext,
    thick,
    minimum size = 1cm,
  },
  every edge/.style = {draw=tikz-edge, thick, ->, >=latex},
}
\end{lstlisting}

\subsection{DFA: autómata determinista}
\label{subsec:dfa}

DFA para el lenguaje $L = \{\eqvec{w} \in \eqdom{\{a,b\}^*} \mid
\eqvec{w}$ contiene el sufijo $\eqk{abb}\}$~\cite{sipser2012introduction}:

\begin{figure}[htbp]
  \centering
  \begin{tikzpicture}[
    shorten >=1pt, node distance=2.5cm, on grid, auto, thick,
    state/.style={circle, draw=tikz-node-border, fill=background,
                  text=maintext, minimum size=1cm},
    every edge/.style={draw=tikz-edge, ->, >=latex},
  ]
    \node[state, initial, initial text={}, draw=tikz-initial]
      (A) {$\eqk{A}$};
    \node[state, right=of A] (B) {$\eqk{B}$};
    \node[state, right=of B] (D) {$\eqk{D}$};
    \node[state, accepting, draw=tikz-accept, right=of D]
      (E) {$\eqk{E}$};

    \path[->]
      (A) edge[loop above] node {$\eqk{b}$} (A)
          edge              node {$\eqk{a}$} (B)
      (B) edge[loop above] node {$\eqk{a}$} (B)
          edge              node {$\eqk{b}$} (D)
      (D) edge              node {$\eqk{b}$} (E)
          edge[bend left]   node {$\eqk{a}$} (B)
      (E) edge              node[swap] {$\eqk{a}$} (B)
          edge[bend right]  node[swap] {$\eqk{b}$} (A);
  \end{tikzpicture}
  \caption{DFA mínimo para $L = \{\eqvec{w} \in \eqdom{\{a,b\}^*} \mid
           \eqvec{w}$ termina en $\eqk{abb}\}$.
           Estado inicial en azul; estado aceptador en verde.}
  \label{fig:dfa-abb}
\end{figure}

\subsection{NFA con transiciones $\varepsilon$}
\label{subsec:nfa}

Construcción de Thompson para $(\eqk{a} \eqop{\mid} \eqk{b})\eqop{*}\eqk{abb}$.
Los estados intermedios usan el color por defecto; el estado aceptador
usa \texttt{tikz-accept}:

\begin{figure}[htbp]
  \centering
  \begin{tikzpicture}[
    shorten >=1pt, node distance=2cm, on grid, auto, thick,
    state/.style={circle, draw=tikz-node-border, fill=background,
                  text=maintext, minimum size=0.9cm},
    every edge/.style={draw=tikz-edge, ->, >=latex},
  ]
    \node[state, initial, initial text={}, draw=tikz-initial]
      (0) {$\eqk{0}$};
    \node[state, right=of 0] (1) {$\eqk{1}$};
    \node[state, right=of 1] (2) {$\eqk{2}$};
    \node[state, accepting, draw=tikz-accept, right=of 2]
      (3) {$\eqk{3}$};

    \path[->]
      (0) edge[loop above] node {$\eqk{a}\eqop{,}\eqk{b}$} (0)
          edge              node {$\eqk{a}$}                 (1)
      (1) edge              node {$\eqk{b}$}                 (2)
      (2) edge              node {$\eqk{b}$}                 (3);
  \end{tikzpicture}
  \caption{NFA simplificado para $({\eqk{a} \eqop{|} \eqk{b}})^* \eqk{abb}$.
           El estado aceptador $\eqk{3}$ en verde distingue el criterio
           de aceptación.}
  \label{fig:nfa-abb}
\end{figure}

%───────────────────────────────────────────────────────────────────────────────
\section{Árboles}
\label{sec:arboles}
%───────────────────────────────────────────────────────────────────────────────

\subsection{Árbol binario de búsqueda con TikZ}
\label{subsec:bst}

Los árboles binarios se construyen con TikZ usando el estilo de nodo
estándar. El nodo raíz puede destacarse con \texttt{tikz-highlight}:

\begin{figure}[htbp]
  \centering
  \begin{tikzpicture}[
    level distance  = 1.5cm,
    level 1/.style  = {sibling distance=4cm},
    level 2/.style  = {sibling distance=2cm},
    level 3/.style  = {sibling distance=1.2cm},
    every node/.style = {
      circle, draw=tikz-node-border, fill=tikz-node-bg,
      text=maintext, minimum size=0.8cm, thick,
    },
    edge from parent/.style = {draw=tikz-edge, thick},
  ]
    \node[fill=tikz-highlight] {$\eqk{50}$}
      child { node {$\eqk{30}$}
        child { node {$\eqk{20}$}
          child { node {$\eqk{10}$} }
          child[missing] {}
        }
        child { node {$\eqk{40}$} }
      }
      child { node {$\eqk{70}$}
        child { node {$\eqk{60}$} }
        child { node {$\eqk{80}$}
          child[missing] {}
          child { node {$\eqk{90}$} }
        }
      };
  \end{tikzpicture}
  \caption{BST con raíz $\eqk{50}$ destacada. La propiedad de orden:
           todo nodo izquierdo $< \eqk{k} <$ nodo derecho~\cite{cormen2022algorithms}.}
  \label{fig:bst}
\end{figure}

\subsection{Árbol sintáctico con forest}
\label{subsec:parse-tree}

El paquete \texttt{forest}~\cite{forest2017} es más expresivo que TikZ
para árboles con muchos niveles. Ideal para árboles de derivación gramatical
y árboles de expresión:

\begin{lstlisting}[style=latex, frame=none, numbers=none]
\begin{forest}
  for tree={
    draw          = tikz-node-border,
    fill          = tikz-node-bg,
    text          = maintext,
    rounded corners = 2pt,
    s sep         = 1.5em,
    l sep         = 2em,
    edge          = {draw=tikz-edge, thick},
  }
  [$\eqop{+}$
    [$\eqop{\times}$
      [$\eqk{3}$]
      [$\eqk{4}$]
    ]
    [$\eqk{5}$]
  ]
\end{forest}
\end{lstlisting}

\begin{observacion}{qtree para árboles simples}
Para árboles lingüísticos o de derivación sencillos, \texttt{qtree} ofrece
una sintaxis más compacta:
\begin{lstlisting}[style=latex, frame=none, numbers=none]
\Tree [.S [.NP [.Det la ] [.N funcion ] ]
          [.VP [.V converge ] [.PP [.P en ] [.NP \(L^2\) ] ] ] ]
\end{lstlisting}
\texttt{forest} es preferible cuando se necesita control fino sobre
colores de nodos individuales o layouts no estándar.
\end{observacion}

%───────────────────────────────────────────────────────────────────────────────
\section{Estructuras de datos}
\label{sec:estructuras-datos}
%───────────────────────────────────────────────────────────────────────────────

Los diagramas de estructuras de datos (arrays, tablas hash, listas enlazadas)
comparten un vocabulario visual: celdas rectangulares para memoria, flechas
para punteros, color para estado semántico (libre, ocupado, colisión).

\subsection{Tabla hash con separate chaining}
\label{subsec:hash-chaining}

El diagrama muestra una tabla hash donde el índice $\eqop{h}(\eqk{k}) =
\eqk{k} \eqop{\bmod} \eqk{m}$ mapea claves a buckets. Los buckets
ocupados usan \texttt{tikz-accept} (Green); los vacíos, \texttt{tikz-node-bg}
(Surface0); las colisiones resueltas por chaining extienden la cadena
hacia la derecha~\cite{cormen2022algorithms}:

\begin{figure}[htbp]
  \centering
  \begin{tikzpicture}[font=\small,
    cell/.style = {
      draw=tikz-node-border, minimum width=1.4cm, minimum height=0.7cm,
    },
    occ/.style  = {cell, fill=tikz-accept,   text=background},
    empty/.style= {cell, fill=tikz-node-bg,  text=maintext},
    key/.style  = {cell, fill=tikz-initial,  text=background},
    arr/.style  = {draw=tikz-edge, ->, >=latex, thick},
  ]
    % --- Buckets (índices 0-7) ---
    \foreach \i in {0,...,7} {
      \node[empty] (b\i) at (0, {-\i*0.9}) {};
      \node[left=0.15cm of b\i, text=maintext] {\(\i\)};
    }

    % --- Buckets ocupados ---
    \node[occ] at (b1) {\(\eqk{21}\)};
    \node[occ] at (b3) {\(\eqk{12}\)};

    % --- Chains desde índice 1: 21→31→41 ---
    \node[occ] (n1a) at (2.2, {-1*0.9}) {\(\eqk{31}\)};
    \node[occ] (n1b) at (4.0, {-1*0.9}) {\(\eqk{41}\)};
    \draw[arr] (b1)  -- (n1a);
    \draw[arr] (n1a) -- (n1b);

    % --- Chain desde índice 3: 12→22 ---
    \node[occ] (n3a) at (2.2, {-3*0.9}) {\(\eqk{22}\)};
    \draw[arr] (b3) -- (n3a);

    % --- Etiquetas ---
    \node[above=0.4cm of b0, text=maintext] {\textbf{Índice}};
    \node[above=0.4cm of n1a, text=maintext] {\textbf{Cadena}};
  \end{tikzpicture}
  \caption{Separate chaining: los buckets ocupados (verde) encadenan
           los elementos con el mismo índice hash. Vacíos en gris.}
  \label{fig:chaining}
\end{figure}

\subsection{Linear probing y colisión}
\label{subsec:linear-probing}

En linear probing, una colisión en el índice $\eqop{h}(\eqk{k})$ activa
una búsqueda lineal del siguiente slot libre. El slot de colisión usa
\texttt{tikz-highlight} para visualizar el conflicto; el slot final usa
\texttt{tikz-accept}:

\begin{figure}[htbp]
  \centering
  \begin{tikzpicture}[font=\small,
    cell/.style  = {draw=tikz-node-border, minimum width=0.9cm,
                    minimum height=0.65cm},
    occ/.style   = {cell, fill=tikz-accept,    text=background},
    coll/.style  = {cell, fill=tikz-highlight, text=maintext},
    empty/.style = {cell, fill=tikz-node-bg,   text=maintext},
    arr/.style   = {draw=tikz-edge, ->, >=latex, thick, dashed},
  ]
    % Array de 10 celdas
    \foreach \i in {0,...,9} {
      \node[empty] (a\i) at ({\i*1.1}, 0) {};
      \node[below=0.15cm of a\i, text=tikz-label, font=\scriptsize] {\(\i\)};
    }

    % Valor 15 en índice 5 (sin colisión)
    \node[occ] at (a5) {\(\eqk{15}\)};

    % Valor 25: colisiona en 5, va a 6
    \node[coll] at (a5) {\(\eqk{15}\)};   % ya ocupado
    \node[occ]  at (a6) {\(\eqk{25}\)};

    % Flecha de probing
    \draw[arr] (a5.north) to[bend left=40]
      node[above, text=maintext, font=\scriptsize] {probing}
      (a6.north);

    % Etiquetas semánticas
    \node[above=1.2cm of a5, text=tikz-highlight]
      {\small colisión: $\eqop{h}(\eqk{25}) = \eqk{5}$};
    \node[above=1.2cm of a6, text=tikz-accept]
      {\small colocado en $\eqk{6}$};
  \end{tikzpicture}
  \caption{Linear probing: $\eqk{25}$ colisiona en índice $\eqop{h}(\eqk{25}) = \eqk{5}$
           y se desplaza al siguiente slot libre $\eqk{6}$.}
  \label{fig:linear-probing}
\end{figure}

\subsection{Diagrama de hashing completo}
\label{subsec:hash-completo}

El diagrama más informativo muestra claves, función hash y buckets
como tres columnas conectadas. La función hash usa \texttt{env-phys}
(Sapphire) por ser una transformación matemática; las claves de entrada
usan \texttt{tikz-initial} (Blue):

\begin{figure}[htbp]
  \centering
  \begin{tikzpicture}[font=\small,
    key/.style   = {draw=tikz-node-border, fill=tikz-initial,
                    text=background, minimum width=2.2cm,
                    minimum height=0.7cm},
    fn/.style    = {draw=tikz-node-border, fill=env-phys,
                    text=background, minimum width=1.3cm,
                    minimum height=3.8cm, align=center},
    occ/.style   = {draw=tikz-node-border, fill=tikz-accept,
                    text=background, minimum width=1.4cm,
                    minimum height=0.55cm},
    empty/.style = {draw=tikz-node-border, fill=tikz-node-bg,
                    text=maintext, minimum width=1.4cm,
                    minimum height=0.55cm},
    arr/.style   = {draw=tikz-edge, ->, >=latex, thick},
  ]
    % --- Claves ---
    \node[key] (k1) at (-4.5,  1.2) {John Smith};
    \node[key] (k2) at (-4.5,  0.0) {Lisa Smith};
    \node[key] (k3) at (-4.5, -1.2) {Sandra Dee};

    % --- Función hash ---
    \node[fn] (hf) at (-2, 0) {$\eqop{h}$\\function};

    % --- Buckets ---
    \foreach \i in {0,...,5} {
      \node[empty] (bk\i) at (1, {1.5 - \i*0.6}) {};
      \node[left=0.1cm of bk\i, text=tikz-label,
            font=\scriptsize] {\(\i\)};
    }
    \node[occ] at (bk1) {\footnotesize 521-8976};
    \node[occ] at (bk2) {\footnotesize 521-1234};
    \node[occ] at (bk4) {\footnotesize 521-9655};

    % --- Aristas claves → hash ---
    \draw[arr] (k1) -- (hf);
    \draw[arr] (k2) -- (hf);
    \draw[arr] (k3) -- (hf);

    % --- Aristas hash → buckets ---
    \draw[arr] (hf.east) to[out=0, in=180] (bk1.west);
    \draw[arr] (hf.east) to[out=0, in=180] (bk2.west);
    \draw[arr] (hf.east) to[out=0, in=180] (bk4.west);

    % --- Etiquetas ---
    \node[above=0.4cm of k1,   text=maintext] {\textbf{Claves}};
    \node[above=0.4cm of hf,   text=maintext] {\textbf{Función hash}};
    \node[above=0.4cm of bk0,  text=maintext] {\textbf{Buckets}};
  \end{tikzpicture}
  \caption{Hashing: claves (azul) transformadas por
           $\eqop{h}$ (sapphire) a buckets (verde=ocupado, gris=libre).
           La función hash tiene rol de transformación, análogo a un
           operador físico.}
  \label{fig:hashing}
\end{figure}

%───────────────────────────────────────────────────────────────────────────────
\section{Grafos}
\label{sec:grafos}
%───────────────────────────────────────────────────────────────────────────────

\subsection{Grafo dirigido ponderado}
\label{subsec:grafo-dirigido}

Los pesos de aristas son datos conocidos del problema — usan
\texttt{eq-known} (Mauve). Los nodos son genéricos salvo el nodo fuente
que usa \texttt{tikz-initial} y el nodo destino \texttt{tikz-accept}:

\begin{figure}[htbp]
  \centering
  \begin{tikzpicture}[
    node distance = 2.5cm,
    vertex/.style = {circle, draw=tikz-node-border, fill=tikz-node-bg,
                     text=maintext, minimum size=0.9cm, thick},
    every edge/.style = {draw=tikz-edge, ->, >=latex, thick},
    wlabel/.style = {text=eq-known, font=\small},
  ]
    \node[vertex, draw=tikz-initial, fill=tikz-initial!15] (s) {$\eqk{s}$};
    \node[vertex, right=of s]  (u) {$\eqk{u}$};
    \node[vertex, right=of u]  (v) {$\eqk{v}$};
    \node[vertex, below=of u]  (x) {$\eqk{x}$};
    \node[vertex, draw=tikz-accept, fill=tikz-accept!15,
          right=of x]          (t) {$\eqk{t}$};

    \path[->]
      (s) edge node[above, wlabel] {$\eqk{6}$} (u)
      (s) edge node[left,  wlabel] {$\eqk{7}$} (x)
      (u) edge node[above, wlabel] {$\eqk{5}$} (v)
      (u) edge node[right, wlabel] {$\eqk{8}$} (x)
      (u) edge node[right, wlabel] {$\eqk{-4}$}(t)
      (x) edge node[right, wlabel] {$\eqk{9}$} (v)
      (x) edge node[below, wlabel] {$\eqk{7}$} (t)
      (v) edge node[right, wlabel] {$\eqk{-2}$}(t);
  \end{tikzpicture}
  \caption{Grafo dirigido ponderado. Fuente $\eqk{s}$ (azul),
           destino $\eqk{t}$ (verde). Pesos en malva
           como datos conocidos del problema de camino mínimo~\cite{cormen2022algorithms}.}
  \label{fig:grafo-ponderado}
\end{figure}

\subsection{Nodo activo durante ejecución de algoritmo}
\label{subsec:grafo-highlight}

Para ilustrar la ejecución paso a paso de Dijkstra u otro algoritmo
sobre un grafo, \texttt{tikz-highlight} marca el nodo activo en cada
paso. El patrón es sobreescribir el estilo del nodo localmente:

\begin{lstlisting}[style=latex, frame=none, numbers=none]
% Nodo activo en el paso actual del algoritmo
\node[vertex, fill=tikz-highlight] (u) {$\eqk{u}$};

% Arista relajada en este paso
\path[->, draw=tikz-highlight, very thick]
  (s) edge node[above, text=tikz-highlight] {$\eqk{6}$} (u);

% Nodo ya resuelto (distancia final)
\node[vertex, fill=tikz-accept!20, draw=tikz-accept] (s) {$\eqk{s}$};
\end{lstlisting}

%───────────────────────────────────────────────────────────────────────────────
\section{Diagramas conmutativos}
\label{sec:tikz-cd}
%───────────────────────────────────────────────────────────────────────────────

\begin{definicion}{tikz-cd}
\label{def:tikzcd}
El paquete \texttt{tikz-cd}~\cite{tikzcd2021} extiende TikZ para producir
diagramas conmutativos con sintaxis matricial. Cada celda es un nodo;
las flechas se declaran con \texttt{\textbackslash arrow[dirección]\{etiqueta\}}.
Las etiquetas de morfismos son objetos matemáticos y aplican el coloreado
semántico de ecuaciones.
\end{definicion}

\subsection{Diagrama conmutativo: teorema del isomorfismo}
\label{subsec:cd-isomorfismo}

\begin{figure}[htbp]
  \centering
  \begin{tikzcd}[
    every arrow/.append style = {draw=tikz-edge, thick},
    every label/.append style = {text=maintext},
  ]
    \eqdom{G}
      \arrow[r, "\eqop{\varphi}"]
      \arrow[d, "\eqop{\pi}"', two heads]
    & \eqdom{H} \\
    \eqdom{G} \eqop{/} \eqdom{\ker\varphi}
      \arrow[ur, dashed, "\eqop{\tilde{\varphi}}"', "\sim"']
  \end{tikzcd}
  \caption{Primer teorema de isomorfismo: $\eqdom{G}/\eqdom{\ker\varphi}
           \cong \eqop{\mathrm{im}}(\eqop{\varphi})$.
           $\eqop{\pi}$ es la proyección canónica (flecha doble);
           $\eqop{\tilde{\varphi}}$ es el isomorfismo inducido (punteado).}
  \label{fig:cd-isomorfismo}
\end{figure}

\subsection{Diagrama conmutativo: sucesión exacta corta}
\label{subsec:cd-ses}

\begin{figure}[htbp]
  \centering
  \begin{tikzcd}[
    every arrow/.append style = {draw=tikz-edge, thick},
  ]
    \eqdom{0}
      \arrow[r]
    & \eqdom{A}
      \arrow[r, "\eqop{f}"]
      \arrow[d, "\eqop{\alpha}"']
    & \eqdom{B}
      \arrow[r, "\eqop{g}"]
      \arrow[d, "\eqop{\beta}"']
    & \eqdom{C}
      \arrow[r]
      \arrow[d, "\eqop{\gamma}"]
    & \eqdom{0} \\
    \eqdom{0}
      \arrow[r]
    & \eqdom{A'}
      \arrow[r, "\eqop{f'}"']
    & \eqdom{B'}
      \arrow[r, "\eqop{g'}"']
    & \eqdom{C'}
      \arrow[r]
    & \eqdom{0}
  \end{tikzcd}
  \caption{Diagrama del lema de los cinco: morfismo entre dos sucesiones
           exactas cortas. La conmutatividad implica que si $\eqop{\alpha}$
           y $\eqop{\gamma}$ son isomorfismos, también lo es
           $\eqop{\beta}$~\cite{mac2013categories}.}
  \label{fig:cd-five-lemma}
\end{figure}

\begin{observacion}{Estilos de flecha en tikz-cd}
Las flechas más comunes en diagramas conmutativos y su declaración:
\begin{lstlisting}[style=latex, frame=none, numbers=none]
\arrow[r]                  % flecha estandar ->
\arrow[r, two heads]       % epimorfico ->>
\arrow[r, hook]            % monomorfico (gancho)
\arrow[r, dashed]          % morfismo inducido
\arrow[r, Rightarrow]      % transformacion natural =>
\arrow[r, no head]         % relacion de equivalencia
\end{lstlisting}
\end{observacion}

%───────────────────────────────────────────────────────────────────────────────
\section{Patrones de código reutilizables}
\label{sec:tikz-patrones}
%───────────────────────────────────────────────────────────────────────────────

El patrón de definición de estilos al inicio del entorno
\texttt{tikzpicture} es el más limpio y portable: evita contaminar el
ámbito global y hace el diagrama autocontenido.

\begin{lstlisting}[style=latex, frame=none, numbers=none]
\begin{tikzpicture}[
  % Definir estilos locales que consumen aliases semanticos
  vertex/.style = {
    circle,
    draw    = tikz-node-border,
    fill    = tikz-node-bg,
    text    = maintext,
    thick,
    minimum size = 1cm,
  },
  edge/.style = {
    draw    = tikz-edge,
    ->,
    >=latex,
    thick,
  },
  % Variantes semanticas por tipo de nodo
  initial/.style  = {vertex, draw=tikz-initial},
  accepted/.style = {vertex, draw=tikz-accept},
  active/.style   = {vertex, fill=tikz-highlight},
  error/.style    = {vertex, fill=fb-error, text=background},
]
  % ... contenido del diagrama
\end{tikzpicture}
\end{lstlisting}

\begin{observacion}{Reutilización entre figuras}
Si varios diagramas de un capítulo comparten el mismo vocabulario visual,
los estilos se pueden declarar una sola vez con
\texttt{\textbackslash tikzset\{...\}} al inicio del capítulo.
Para estilos que aplican a todo el documento (como el estilo base de
autómatas), declararlos en \texttt{config/} o en un bloque de
\texttt{main.tex} después de cargar los paquetes.
\end{observacion}

%───────────────────────────────────────────────────────────────────────────────
\printbibliography[heading=subbibliography]
